%% SYNTHETIC TEST UNIT generated by writing/tools/make_outline_book.py from FINAL_OUTLINE_v2 -- not book text
\chapter{QCD 基础——从规范不变性到渐近自由}\label{ch:qcd}
本单元为测试用的合成文本（由 FINAL\_OUTLINE\_v2 自动生成），陈述内容从略。

\section{第 1 节（测试）}\label{sec:qcd-why}

\begin{definition}[规范群 SU($N$) 及其表示]\label{def:qcd-gauge-group}
（测试用合成陈述：内容从略。）
% uses: none
\end{definition}

\begin{lemma}[色代数恒等式]\label{lem:qcd-colour-identities}
（测试用合成陈述：内容从略。）
\uses{def:qcd-gauge-group}
\end{lemma}

\begin{example}[单圈图的色因子与双线记号]\label{ex:qcd-colour-algebra}
（测试用合成陈述：内容从略。）
\uses{lem:qcd-colour-identities}
\end{example}

\begin{proposition}[协变导数与场强的规范协变性 \tagR]\label{prop:qcd-covariance}
（$R_{\mathrm{def}}$：由算符定义直接推出。）
（测试用合成陈述：内容从略。）
（本命题只涉及部分子态。）
\uses{def:qcd-gauge-group, lem:qcd-colour-identities}
\end{proposition}

\section{第 2 节（测试）}\label{sec:qcd-gauge}

\begin{definition}[QCD 拉氏量与规范变换]\label{def:prelim-qcd}
（测试用合成陈述：内容从略。）
\uses{def:qcd-gauge-group}
\end{definition}

\begin{theorem}[经典 QCD 拉氏量的规范不变性与唯一性 \tagR]\label{thm:qcd-gauge-invariance}
（$R_{\mathrm{def}}$：由算符定义直接推出。）
（测试用合成陈述：内容从略。）
\uses{def:prelim-qcd, prop:qcd-covariance}
\end{theorem}

\begin{remark}[$\theta$ 项]\label{rem:qcd-theta}
（测试用合成陈述：内容从略。）
% uses: none
\end{remark}

\begin{remark}[Faddeev–Popov 构造（形式推导）]\label{rem:qcd-fp-derivation}
（测试用合成陈述：内容从略。）
% uses: none
\end{remark}

\section{第 3 节（测试）}\label{sec:prelim-qcd}

\begin{definition}[协变规范、Faddeev–Popov 鬼场与微扰 QCD]\label{def:qcd-fp}
（测试用合成陈述：内容从略。）
\uses{def:prelim-qcd}
\end{definition}

\begin{example}[由拉氏量导出传播子、鬼场顶点与闭合鬼圈因子]\label{ex:qcd-fp-lowest-order}
（测试用合成陈述：内容从略。）
\uses{def:qcd-fp}
\end{example}

\begin{definition}[Feynman 规则]\label{def:prelim-feynman-rules}
（测试用合成陈述：内容从略。）
\uses{def:prelim-qcd, def:qcd-fp}
\end{definition}

\begin{example}[树图 $q\bar q\to gg$：为什么需要鬼场]\label{ex:qcd-qqbar-gg}
（测试用合成陈述：内容从略。）
\uses{def:prelim-feynman-rules}
\end{example}

\section{第 4 节（测试）}\label{sec:qcd-quantization}

\begin{definition}[维数正规化]\label{def:qcd-dimreg}
（测试用合成陈述：内容从略。）
% uses: none
\end{definition}

\begin{proposition}[紫外超表观发散度 \tagR]\label{prop:qcd-superficial-degree}
（$R_{\mathrm{def}}$：由算符定义直接推出。）
（测试用合成陈述：内容从略。）
\uses{def:prelim-feynman-rules}
\end{proposition}

\begin{definition}[BRST 变换]\label{def:prelim-brst}
（测试用合成陈述：内容从略。）
\uses{def:qcd-fp}
\end{definition}

\begin{lemma}[$s^2=0$ 与拉氏量的 BRST 不变性 \tagR]\label{lem:prelim-brst-nilpotent}
（$R_{\mathrm{def}}$：由算符定义直接推出。）
（测试用合成陈述：内容从略。）
\uses{def:prelim-brst, def:prelim-qcd, def:qcd-fp, lem:qcd-colour-identities, prop:qcd-covariance}
\end{lemma}

\section{第 5 节（测试）}\label{sec:qcd-regularization}

\begin{proposition}[自由渐近 Fock 空间上的 BRST 上同调（Kugo–Ojima 四重态） \tagR]\label{prop:qcd-brst-cohomology}
（$R_{\mathrm{def}}$：由算符定义直接推出。）
（测试用合成陈述：内容从略。）
\uses{def:prelim-brst, def:qcd-fp, lem:prelim-brst-nilpotent}
\end{proposition}

\begin{example}[单一动量模式的四重态]\label{ex:qcd-brst-quartet}
（测试用合成陈述：内容从略。）
\uses{prop:qcd-brst-cohomology, ex:qcd-qqbar-gg}
\end{example}

\begin{lemma}[维数正规化理论的 Slavnov–Taylor 恒等式 \tagR]\label{lem:qcd-st-regularized}
（$R_{\mathrm{PT}}$：在重整化微扰论中对红外正规化的部分子态严格。）
（测试用合成陈述：内容从略。）
\uses{def:qcd-dimreg, lem:prelim-brst-nilpotent, def:prelim-feynman-rules}
\end{lemma}

\begin{definition}[重整化常数与抵消项]\label{def:qcd-renormalization}
（测试用合成陈述：内容从略。）
\uses{def:prelim-qcd, def:qcd-fp, def:qcd-dimreg}
\end{definition}

\section{第 6 节（测试）}\label{sec:prelim-brst}

\begin{definition}[$\overline{\rm MS}$ 方案]\label{def:prelim-msbar}
（测试用合成陈述：内容从略。）
\uses{def:qcd-renormalization}
\end{definition}

\begin{theorem}[QCD 的可重整性 \tagR]\label{thm:qcd-renormalizability}
（$R_{\mathrm{PT}}$：在重整化微扰论中对红外正规化的部分子态严格。）
（测试用合成陈述：内容从略。）
（本命题只涉及部分子态。）
\uses{prop:qcd-superficial-degree, def:qcd-renormalization, def:prelim-msbar, lem:qcd-st-regularized, thm:qcd-gauge-invariance}
\end{theorem}

\begin{proposition}[重整化常数间的 Slavnov–Taylor 关系 \tagR]\label{prop:qcd-st-relations}
（$R_{\mathrm{PT}}$：在重整化微扰论中对红外正规化的部分子态严格。）
（测试用合成陈述：内容从略。）
（本命题只涉及部分子态。）
\uses{thm:qcd-renormalizability, lem:qcd-st-regularized, def:qcd-renormalization}
\end{proposition}

\begin{theorem}[Slavnov–Taylor 恒等式（因子化所用形式） \tagR]\label{thm:prelim-slavnov-taylor}
（$R_{\mathrm{PT}}$：在重整化微扰论中对红外正规化的部分子态严格。）
（测试用合成陈述：内容从略。）
（本命题只涉及部分子态。）
\uses{lem:qcd-st-regularized, thm:qcd-renormalizability, def:prelim-brst}
\end{theorem}

\section{第 7 节（测试）}\label{sec:prelim-renorm}

\begin{proposition}[物理量的规范参数无关性 \tagR]\label{prop:qcd-xi-independence}
（$R_{\mathrm{PT}}$：在重整化微扰论中对红外正规化的部分子态严格。）
（测试用合成陈述：内容从略。）
（本命题只涉及部分子态。）
\uses{thm:prelim-slavnov-taylor, thm:qcd-renormalizability, lem:prelim-brst-nilpotent, prop:qcd-brst-cohomology}
\end{proposition}

\begin{theorem}[规范不变复合算符的重整化（Joglekar–Lee） \tagR]\label{thm:prelim-composite-renorm}
（$R_{\mathrm{PT}}$：在重整化微扰论中对红外正规化的部分子态严格。）
（测试用合成陈述：内容从略。）
\uses{thm:qcd-renormalizability, lem:qcd-st-regularized, thm:qcd-gauge-invariance}
\end{theorem}

\begin{remark}[BPHZ 减除与 Slavnov–Taylor 恒等式；减除度的顶点可加性]\label{rem:qcd-bphz-brst}
（测试用合成陈述：内容从略。）
% uses: none
\end{remark}

\begin{example}[单圈胶子自能与 $Z_3$]\label{ex:qcd-gluon-self-energy}
（测试用合成陈述：内容从略。）
\uses{def:prelim-feynman-rules, def:qcd-renormalization, def:prelim-msbar, lem:qcd-colour-identities, lem:qcd-st-regularized}
\end{example}

\section{第 8 节（测试）}\label{sec:qcd-renorm}

\begin{example}[夸克自能与夸克—胶子顶点]\label{ex:qcd-quark-self-energy-vertex}
（测试用合成陈述：内容从略。）
\uses{def:prelim-feynman-rules, def:qcd-renormalization, ex:qcd-colour-algebra}
\end{example}

\begin{example}[用四种 Slavnov–Taylor 等价途径计算 $\beta_0$]\label{ex:qcd-beta0-four-ways}
（测试用合成陈述：内容从略。）
\uses{prop:qcd-st-relations, ex:qcd-gluon-self-energy, ex:qcd-quark-self-energy-vertex}
\end{example}

\begin{proposition}[$\beta$ 函数 \tagR]\label{prop:prelim-beta}
（$R_{\mathrm{PT}}$：在重整化微扰论中对红外正规化的部分子态严格。）
（测试用合成陈述：内容从略。）
\uses{prop:qcd-st-relations, def:prelim-msbar, ex:qcd-beta0-four-ways}
\end{proposition}

\begin{proposition}[重整化群方程 \tagR]\label{prop:qcd-rge}
（$R_{\mathrm{PT}}$：在重整化微扰论中对红外正规化的部分子态严格。）
（测试用合成陈述：内容从略。）
（本命题只涉及部分子态。）
\uses{thm:qcd-renormalizability, def:prelim-msbar, prop:prelim-beta, prop:qcd-xi-independence}
\end{proposition}

\section{第 9 节（测试）}\label{sec:qcd-running}

\begin{proposition}[重整化群改进 \tagR]\label{prop:prelim-rg}
（$R_{\mathrm{PT}}$：在重整化微扰论中对红外正规化的部分子态严格。）
（测试用合成陈述：内容从略。）
\uses{prop:qcd-rge, prop:prelim-beta}
\end{proposition}

\begin{definition}[$\Lambda_{\overline{\rm MS}}$ 与跑动耦合]\label{def:qcd-lambda}
（测试用合成陈述：内容从略。）
\uses{prop:prelim-beta}
\end{definition}

\begin{remark}[重夸克退耦与有效味数]\label{rem:qcd-decoupling}
（测试用合成陈述：内容从略。）
% uses: none
\end{remark}

\begin{example}[从 $\alpha_s(M_Z)$ 出发的数值跑动]\label{ex:qcd-running-numbers}
（测试用合成陈述：内容从略。）
\uses{def:qcd-lambda, prop:prelim-rg}
\end{example}

\section{第 10 节（测试）}\label{sec:qcd-af}

\begin{theorem}[渐近自由 \tagR]\label{thm:qcd-asymptotic-freedom}
（$R_{\mathrm{PT}}$：在重整化微扰论中对红外正规化的部分子态严格。）
（测试用合成陈述：内容从略。）
\uses{prop:prelim-beta, prop:qcd-rge, prop:qcd-xi-independence}
\end{theorem}

\begin{remark}[屏蔽与反屏蔽的直观图像]\label{rem:qcd-antiscreening}
（测试用合成陈述：内容从略。）
% uses: none
\end{remark}

\begin{remark}[渐近自由意味着什么、不意味着什么]\label{rem:qcd-af-meaning}
（测试用合成陈述：内容从略。）
% uses: none
\end{remark}

\begin{proposition}[树图层次的软极限与共线极限 \tagR]\label{prop:qcd-soft-collinear-limits}
（$R_{\mathrm{PT}}$：在重整化微扰论中对红外正规化的部分子态严格。）
（测试用合成陈述：内容从略。）
（本命题只涉及部分子态。）
\uses{def:prelim-feynman-rules, lem:qcd-colour-identities}
\end{proposition}

\section{第 11 节（测试）}\label{sec:qcd-ir}

\begin{definition}[红外安全（非形式）]\label{def:qcd-ir-safety-informal}
（测试用合成陈述：内容从略。）
% uses: none
\end{definition}

\begin{example}[$e^+e^-\to q\bar q(g)$：实辐射与虚修正的相消]\label{ex:qcd-ee-real-virtual}
（测试用合成陈述：内容从略。）
\uses{prop:qcd-soft-collinear-limits, def:qcd-ir-safety-informal, def:qcd-dimreg}
\end{example}

\begin{example}[第一个红外安全的排他量：Sterman–Weinberg 两喷注截面]\label{ex:qcd-sterman-weinberg}
（测试用合成陈述：内容从略。）
\uses{ex:qcd-ee-real-virtual, def:qcd-ir-safety-informal}
\end{example}

\begin{remark}[红外定理预览]\label{rem:qcd-ir-theorems-preview}
（测试用合成陈述：内容从略。）
% uses: none
\end{remark}
